\documentclass[aspectratio=169]{beamer}
\usetheme{Madrid}
\usecolortheme{default}
\usepackage[utf8]{inputenc}
\usepackage{amsmath}
\usepackage{amsfonts}
\usepackage{amssymb}

\setbeamertemplate{navigation symbols}{}
\setbeamerfont{block body}{size=\small}
\setbeamerfont{alertblock body}{size=\small}

\title{CRISC Domain B -- Impossible Questions}
\subtitle{IT Risk Assessment -- Expert Level Only (Merged Edition)}
\author{CRISC Mastery Series}
\date{\today}

% ------------------------------------------------------------
% \hq{number}{topic}{question}{options}{letter}{explanation}{trap}
% Creates one question frame and one explanation frame.
% (No blank lines are allowed inside the arguments.)
% ------------------------------------------------------------
\newcommand{\hq}[7]{%
	\begin{frame}{Q#1: #2}
		\begin{block}{Question}
			#3
		\end{block}
		\begin{enumerate}[(A)]
			#4
		\end{enumerate}
	\end{frame}
	\begin{frame}{Q#1: Explanation}
		\begin{block}{Correct Answer: #5}
			\begin{itemize}
				#6
			\end{itemize}
		\end{block}
		\begin{alertblock}{Critical Trap}
			#7
		\end{alertblock}
	\end{frame}
}

\begin{document}
	
	\begin{frame}
		\titlepage
	\end{frame}
	
	\begin{frame}{Warning}
		\begin{alertblock}{Extreme Difficulty}
			Every easy and purely definitional question has been removed.\\
			What remains needs calculation, probability, dependency analysis and judgment.\\
			Do not be discouraged if you miss most of them.
		\end{alertblock}
	\end{frame}
	
	\begin{frame}{Outline}
		\tableofcontents
	\end{frame}
	
	% ============================================================
	\section{Quantitative Risk Analysis}
	% ============================================================
	
	\hq{1}{Control ROI}
	{An asset is worth \$2M with an exposure factor of 40\% and an ARO of 0.5. A control reduces the ARO to 0.1 and the exposure factor to 25\% at a cost of \$60K per year. What is the net annual benefit?}
	{\item \$350K
		\item \$410K
		\item \$290K
		\item \$50K}
	{C}
	{\item Before: ALE $=$ \$2M $\times$ 40\% $\times$ 0.5 $=$ \$400K
		\item After: ALE $=$ \$2M $\times$ 25\% $\times$ 0.1 $=$ \$50K
		\item Gross reduction $=$ \$350K
		\item Net benefit $=$ \$350K $-$ \$60K $=$ \$290K}
	{\$350K ignores the control cost. \$50K is the residual ALE, not a benefit.}
	
	\hq{2}{Joint Probability Bounds}
	{Risk A has a 20\% likelihood and Risk B has a 30\% likelihood. They are known to be positively correlated. Which statement MUST be true about the probability that both occur?}
	{\item It equals exactly 6\%
		\item It is at least 50\%
		\item It can exceed 30\%
		\item It is strictly greater than 6\% and at most 20\%}
	{D}
	{\item Independence gives $0.2 \times 0.3 = 6\%$
		\item Positive correlation pushes the joint probability above 6\%
		\item A joint event can never be more likely than its rarer component
		\item Upper bound $= \min(20\%, 30\%) = 20\%$}
	{Correlation changes the joint probability, but it is always capped by the least likely event.}
	
	\hq{3}{Annualizing a Monthly Probability}
	{A threat event occurs in any given month with probability 5\%, independently of other months. What is the probability of at least one occurrence in 12 months?}
	{\item About 46\%
		\item About 54\%
		\item 5\%
		\item 60\%}
	{A}
	{\item $P(\text{none in 12 months}) = 0.95^{12} \approx 0.54$
		\item $P(\text{at least one}) = 1 - 0.54 \approx 0.46$
		\item 54\% is the probability of NO event
		\item $12 \times 5\% = 60\%$ ignores overlap of outcomes}
	{Probabilities of ``at least one'' do not add linearly. Use the complement.}
	
	\hq{4}{Alert Reliability}
	{The probability of a targeted attack this quarter is 10\%. When an attack is under way, the detection alert fires 90\% of the time. When no attack is under way, it fires falsely 20\% of the time. The alert has just fired. What is the probability that an attack is actually under way?}
	{\item 90\%
		\item 10\%
		\item About 33\%
		\item 50\%}
	{C}
	{\item $P(\text{alert}) = 0.9 \times 0.1 + 0.2 \times 0.9 = 0.27$
		\item $P(\text{attack} \mid \text{alert}) = 0.09 / 0.27 = 1/3$
		\item A low base rate makes most alerts false positives
		\item The alert raises the estimate from 10\% to 33\%, not to 90\%}
	{Detection accuracy is not the same as the probability that an alert is real. Base rates matter.}
	
	\hq{5}{Three-Point Estimate}
	{Analysts estimate a loss at \$100K (optimistic), \$200K (most likely) and \$700K (pessimistic). What is the PERT expected loss?}
	{\item \$333K
		\item \$200K
		\item \$267K
		\item \$400K}
	{C}
	{\item PERT mean $= (O + 4M + P)/6$
		\item $= (100 + 800 + 700)/6 \approx \$267K$
		\item \$333K is the simple average of the three values
		\item \$200K is the mode and ignores the long tail}
	{The most likely value is not the expected value when the distribution is skewed.}
	
	\hq{6}{Skewed Confidence Interval}
	{A FAIR analysis gives a 90\% confidence interval of \$200K to \$2M with a most likely value of \$500K. Which interpretation is BEST?}
	{\item Expected annual loss likely exceeds \$500K because of right skew
		\item Expected annual loss is exactly \$500K
		\item The interval is invalid because it is asymmetric
		\item The loss will certainly fall inside the interval}
	{A}
	{\item Loss distributions are typically right-skewed (lognormal-like)
		\item The long upper tail pulls the mean above the mode
		\item Budgeting on the mode understates expected loss
		\item A 90\% interval leaves a 10\% chance of falling outside it}
	{Most likely $\neq$ expected value. Asymmetric intervals are normal for losses.}
	
	\hq{7}{Aggregating Volatility}
	{Two risks each have an annual loss standard deviation of \$300K. Their correlation is 0.5. What is the standard deviation of the combined loss?}
	{\item \$424K
		\item \$300K
		\item \$600K
		\item \$520K}
	{D}
	{\item $\text{Var} = \sigma^2 + \sigma^2 + 2\rho\sigma^2 = 3\sigma^2$
		\item $\text{SD} = 300 \times \sqrt{3} \approx \$520K$
		\item \$424K assumes independence ($\rho = 0$)
		\item \$600K assumes perfect correlation ($\rho = 1$)}
	{Standard deviations do not add. Correlation controls how much diversification you get.}
	
	\hq{8}{Common-Cause Failure}
	{Two preventive controls each stop an attack 90\% of the time. A design review calls them ``independent layers,'' implying a 1\% chance that both fail. Testing reveals a shared dependency (the same identity provider) whose failure disables both controls 4\% of the time. Otherwise they fail independently. What is the probability that both fail?}
	{\item 1\%
		\item 4\%
		\item About 5\%
		\item 10\%}
	{C}
	{\item $P(\text{both fail}) = 0.04 + 0.96 \times 0.10 \times 0.10$
		\item $= 0.04 + 0.0096 \approx 4.96\%$
		\item Residual risk is about five times the ``independent'' estimate
		\item Defense in depth fails when layers share a failure mode}
	{Layered controls are only independent if they share no common dependency.}
	
	\hq{9}{Rule of Three}
	{A control is tested 30 times with zero failures. Which statement about its true failure rate is MOST defensible?}
	{\item It is 0\%
		\item It is at most 3\%
		\item It is at most 30\%
		\item It is at most roughly 10\% with 95\% confidence}
	{D}
	{\item With zero failures in $n$ trials, the 95\% upper bound is about $3/n$
		\item $3/30 = 10\%$
		\item Zero observed failures does not prove zero failures
		\item Small samples leave wide uncertainty}
	{``No failures observed'' is not ``no failures exist.''}
	
	\hq{10}{FAIR Decomposition}
	{Threat event frequency is 20 per year and vulnerability (probability a threat event becomes a loss event) is 15\%. Primary loss per event is \$40K. There is a 30\% chance of a secondary loss of \$200K per event. What is the annualized loss exposure?}
	{\item \$120K
		\item \$720K
		\item \$300K
		\item \$2.0M}
	{C}
	{\item Loss event frequency $= 20 \times 15\% = 3$ per year
		\item Expected loss per event $= 40K + 0.3 \times 200K = \$100K$
		\item ALE $= 3 \times \$100K = \$300K$
		\item \$120K ignores secondary loss; \$720K treats it as certain; \$2.0M ignores vulnerability}
	{Use loss event frequency, not threat event frequency, and weight secondary loss by its probability.}
	
	% ============================================================
	\section{Risk Evaluation}
	% ============================================================
	
	\hq{11}{Ordinal Matrix Flaw}
	{An organization multiplies 1--5 likelihood and impact scores in a 5x5 matrix. A risk scored 5x1 and a risk scored 1x5 both receive 5. What is the MOST fundamental flaw in this approach?}
	{\item Multiplication always overstates risk
		\item Ordinal ranks are not ratio-scale values, so products can misrank risks with very different expected losses
		\item A 5x5 matrix is too small
		\item Scores above 20 cannot be compared}
	{B}
	{\item Scale points are ordered labels, not measured quantities
		\item Equal products can hide very different loss profiles
		\item Products of ordinal ranks can reverse true priority
		\item Calibrated quantitative ranges are more defensible}
	{A numeric-looking score does not make an ordinal scale quantitative.}
	
	\hq{12}{Aggregate Appetite and Correlation}
	{Three risks each have an expected annual loss of \$400K. Aggregate appetite is \$1M and aggregate tolerance is \$1.3M. The risks are believed to be strongly positively correlated. Which position is MOST accurate?}
	{\item Above appetite but within tolerance on an expected basis; correlation does not change the expected sum but fattens the tail, so the loss distribution must be examined
		\item Within appetite because every individual risk is below \$1M
		\item Beyond tolerance because correlated risks add up to more than their expected values
		\item Within appetite because diversification lowers the expected aggregate loss}
	{A}
	{\item Expected losses add: $3 \times \$400K = \$1.2M$
		\item \$1.2M is above appetite (\$1M) and within tolerance (\$1.3M)
		\item Correlation does not change the expected sum
		\item Correlation does increase tail (extreme-loss) exposure}
	{Correlation affects the spread of outcomes, not the sum of expected values.}
	
	\hq{13}{Cost Exceeds Benefit}
	{A control costing \$400K per year would reduce expected loss by only \$150K per year. Without it, residual risk exceeds appetite but is well within capacity. What should the risk practitioner recommend?}
	{\item Implement the control because the risk exceeds appetite
		\item Escalate for formal, documented acceptance by the owner with authority to exceed appetite, while exploring cheaper alternatives
		\item Accept informally to avoid delay
		\item Buy full insurance without further analysis}
	{B}
	{\item Treatment must be economically justified
		\item Exceeding appetite requires explicit authority to accept
		\item Acceptance must be documented and reviewed
		\item Cheaper responses should still be explored}
	{Exceeding appetite does not force a treatment. It forces a decision by the right authority.}
	
	\hq{14}{Mitigation That Backfires}
	{After a new control is deployed, the verified residual risk is HIGHER than the original inherent risk. What is the MOST likely explanation?}
	{\item An arithmetic error is the only plausible cause
		\item The control introduced secondary risk that outweighs its benefit, so the response must be redesigned
		\item Inherent risk was necessarily overstated
		\item The risk has been transferred to the control owner}
	{B}
	{\item Controls can create new risks (complexity, new attack surface, new dependencies)
		\item Secondary risk must be assessed alongside residual risk
		\item Net exposure, not control count, is what matters
		\item Redesign or replace the response}
	{A control is not automatically net-positive.}
	
	% ============================================================
	\section{Business Impact Analysis}
	% ============================================================
	
	\hq{15}{Recovery Dependencies}
	{System A has an RTO of 4 hours. A's own restoration takes 2 hours and can start only after both System B and System C are available. B takes 3 hours to restore and C takes 2 hours, restored in parallel. Can A meet its RTO?}
	{\item Yes, the longest single task is 3 hours
		\item Yes, B and C run in parallel and fit inside 4 hours
		\item No, recovery takes 7 hours
		\item No, recovery takes 5 hours}
	{D}
	{\item B and C run in parallel, so the dependency stage takes $\max(3, 2) = 3$ hours
		\item A then needs its own 2 hours: $3 + 2 = 5$ hours
		\item 5 hours exceeds the 4-hour RTO
		\item 7 hours would be the fully sequential case}
	{Critical path = longest dependency chain, not the longest single task.}
	
	\hq{16}{Probability of Meeting RPO}
	{The RPO is 1 hour and backups run every 4 hours. A failure occurs at a uniformly random time. What is the probability that the data loss stays within the RPO?}
	{\item 0\%
		\item 25\%
		\item 50\%
		\item 75\%}
	{B}
	{\item Data loss equals the time since the last backup, uniform from 0 to 4 hours
		\item Loss is within RPO only when it is at most 1 hour
		\item Probability $= 1/4 = 25\%$
		\item Expected loss is 2 hours, twice the RPO}
	{A backup interval longer than the RPO means the RPO is usually missed.}
	
	\hq{17}{MTD Consistency}
	{A process has an RTO of 4 hours, a work recovery time (WRT) of 6 hours and an MTD of 8 hours. What is the PRIMARY problem?}
	{\item Consistent because RTO (4 h) is less than MTD (8 h)
		\item Consistent because WRT is excluded from MTD
		\item MTD should automatically be raised to 20 hours
		\item Inconsistent: RTO plus WRT is 10 hours, which exceeds the 8-hour MTD}
	{D}
	{\item The business must be fully operating again within the MTD
		\item Requirement: $\text{RTO} + \text{WRT} \leq \text{MTD}$
		\item Here $4 + 6 = 10 > 8$
		\item Shorten RTO or WRT, or revalidate the MTD with the business}
	{Meeting the RTO alone does not guarantee meeting the MTD.}
	
	\hq{18}{MTD from Escalating Impact}
	{Outage losses are \$10K per hour for the first 4 hours and \$50K per hour afterwards. Management defines MTD as the point where cumulative loss reaches \$400K. What is the MTD?}
	{\item 11.2 hours
		\item 8 hours
		\item 4 hours
		\item 40 hours}
	{A}
	{\item First 4 hours: $4 \times \$10K = \$40K$
		\item Remaining threshold: $\$400K - \$40K = \$360K$
		\item At \$50K per hour: $360 / 50 = 7.2$ hours
		\item MTD $= 4 + 7.2 = 11.2$ hours}
	{Impact is not linear. Piece the loss curve together before dividing.}
	
	% ============================================================
	\section{Risk Treatment}
	% ============================================================
	
	\hq{19}{Treatment Selection}
	{A risk has an 80\% likelihood and a \$100K impact. Treatment A costs \$50K and cuts likelihood to 20\%. Treatment B costs \$30K and cuts impact to \$40K. Applying both costs \$80K. Which option has the best net benefit?}
	{\item Both treatments
		\item Treatment B only
		\item Neither
		\item Treatment A only}
	{B}
	{\item Baseline expected loss $= 0.8 \times 100 = \$80K$
		\item A: loss \$20K, net benefit $= 60 - 50 = +\$10K$
		\item B: loss \$32K, net benefit $= 48 - 30 = +\$18K$
		\item Both: loss \$8K, net benefit $= 72 - 80 = -\$8K$}
	{Benefits of overlapping treatments are not additive. The second one has diminishing value.}
	
	\hq{20}{Present Value of a Control}
	{A control costs \$120K upfront and reduces expected annual loss from \$90K to \$20K for 3 years. The discount rate is 10\%. What is the approximate NPV?}
	{\item About $+$\$54K, so approve
		\item About $+$\$90K
		\item About $-$\$50K, so reject
		\item About $+$\$210K}
	{A}
	{\item Annual benefit $= \$70K$
		\item PV factor for 3 years at 10\% $\approx 2.487$
		\item PV of benefits $\approx \$174K$
		\item NPV $\approx 174 - 120 = +\$54K$}
	{\$90K is the undiscounted net. \$210K is gross undiscounted benefit.}
	
	\hq{21}{Insurance Economics}
	{There is a 10\% annual chance of a \$2M loss. Insurance covers 80\% of the loss above a \$200K deductible and costs \$60K per year. What is the total expected annual cost with insurance?}
	{\item \$56K
		\item \$200K
		\item \$116K
		\item \$260K}
	{C}
	{\item Retained per event $= 200K + 20\% \times 1.8M = \$560K$
		\item Expected retained loss $= 0.1 \times 560K = \$56K$
		\item Add premium: $56K + 60K = \$116K$
		\item Uninsured expected loss would be \$200K}
	{Include the premium and the retained share. Insurance also leaves counterparty and coverage risk.}
	
	% ============================================================
	\section{Risk Monitoring}
	% ============================================================
	
	\hq{22}{Control Test Results}
	{A control runs 250 times per year. A sample of 25 executions shows 5 deviations. The tolerable deviation rate is 5\%. What is the BEST conclusion?}
	{\item Effective, because the sample is only 10\% of the population
		\item Effective, because 80\% of items passed
		\item Inconclusive until all 250 executions are tested
		\item Not effective: the 20\% sample deviation rate far exceeds the 5\% tolerable rate, so remediate and expand testing}
	{D}
	{\item Sample deviation rate $= 5/25 = 20\%$
		\item 20\% is four times the tolerable rate
		\item A passing rate of 80\% is not acceptable here
		\item Full-population testing is not needed to see a clear failure}
	{Compare results to the tolerable rate, not to a ``mostly passed'' feeling.}
	
	\hq{23}{Silent Indicators}
	{Every KRI has stayed green for three consecutive years, yet the organization suffered two significant incidents. What is the MOST appropriate conclusion?}
	{\item Controls are exceptionally effective
		\item Thresholds are likely too loose or indicators insensitive, so back-test and recalibrate
		\item Reduce monitoring frequency to save cost
		\item Retire the KRIs}
	{B}
	{\item Indicators should have moved before the incidents
		\item Back-test KRIs against past incidents
		\item Tighten thresholds and add leading indicators
		\item Consistent green is a warning sign, not proof of safety}
	{A KRI that cannot detect a known failure is not doing its job.}
	
	% ============================================================
	\section{Answer Key and Summary}
	% ============================================================
	
	\begin{frame}{Answer Key}
		\begin{columns}[T]
			\begin{column}{0.32\textwidth}
				Q1: C\\ Q2: D\\ Q3: A\\ Q4: C\\ Q5: C\\ Q6: A\\ Q7: D\\ Q8: C
			\end{column}
			\begin{column}{0.32\textwidth}
				Q9: D\\ Q10: C\\ Q11: B\\ Q12: A\\ Q13: B\\ Q14: B\\ Q15: D\\ Q16: B
			\end{column}
			\begin{column}{0.32\textwidth}
				Q17: D\\ Q18: A\\ Q19: B\\ Q20: A\\ Q21: C\\ Q22: D\\ Q23: B
			\end{column}
		\end{columns}
	\end{frame}
	
	\begin{frame}{Key Takeaways for Domain B}
		\begin{block}{Summary}
			\begin{itemize}
				\item \textbf{Analysis:} Use complements, base rates, correlation and skew. Do not add intuitively.
				\item \textbf{Evaluation:} Ordinal scores mislead. Aggregate and tail views matter.
				\item \textbf{BIA:} Critical path, RPO versus backup interval, and RTO $+$ WRT $\leq$ MTD.
				\item \textbf{Treatment:} Compare net, discounted benefit. Overlapping controls are not additive.
				\item \textbf{Monitoring:} Compare to tolerable rates. Silent indicators need back-testing.
			\end{itemize}
		\end{block}
	\end{frame}
	
	\begin{frame}{Thank You}
		\begin{center}
			{\Huge \textbf{Good Luck on Your CRISC Exam!}}
			
			\vspace{1cm}
			
			Remember: the ISACA mindset is about\\
			understanding why a question is asked,\\
			not just what is technically correct.
		\end{center}
	\end{frame}
	
\end{document}